\documentclass{../ElectromagneticFieldTemplate}

% ====== 图片路径设置 ======
\graphicspath{{./Figures/}{../TemplateFigures/}}

% ====== 填写实验基本信息 ======
\setExperimentInfo
    {圆极化波特性的研究}
    {2024217802}
    {严子竣}
    {2024212622}
    {张璇}
    {\today} 

\begin{document}

\maketitle

\section{实验目的}

\begin{enumerate}
    \item 研究右旋、左旋圆极化波的形成、辐射和接收的过程。
    \item 研究右旋、左旋圆极化波的反射和折射特性，及其测试方法。
\end{enumerate}

\section{实验原理与说明}

圆极化波是通过介质片圆波导产生的。当 \(\mathrm{TE}_{10}\) 波经矩—圆波导转换为 \(\mathrm{TE}_{11}\) 波后，在充有介质片的圆波导中分解为垂直于介质片平面的分量 \(E_n\) 和平行于介质片平面的分量 \(E_t\)。由于介质片的存在，两个分量的传播速度不同，选择合适的介质片长度 \(L\) 可使二者产生 \(90^\circ\) 的相位差；适当调整介质片的转角使两分量幅度相等，即可获得圆极化波。

接收圆极化波时，接收端的介质片圆波导同样将入射波分解为两个正交分量。当接收端介质片法线 \(n'\) 与发射端介质片法线 \(n\) 垂直（\(n' \perp n\)）时，可实现同极化接收；当 \(n' \parallel n\) 时，则接收不到同旋向的圆极化波。

圆极化波在介质表面的反射与折射具有旋向反转的特性：右旋圆极化波投射到介质表面时，反射波变为左旋圆极化波，折射波仍为右旋圆极化波；左旋圆极化波投射时，反射波变为右旋圆极化波，折射波仍为左旋圆极化波。入射波、反射波与折射波之间满足能量守恒关系。

\section{实验设备}

电磁波综合测试仪、铝板1块、玻璃板1块、圆喇叭。

\section{实验内容}

\begin{enumerate}
    \item 圆极化波椭圆度的调整与测试。
    \item 圆极化波反射、折射特性的测试。
\end{enumerate}

\section{实验步骤}

\subsection{圆极化波椭圆度的调整与测试}
\begin{enumerate}
    \item 收发端先用矩形喇叭进行对准，确保系统正常工作。
    \item 固定发端电路板、接收端矩形喇叭及接收模块不变。将发端矩形喇叭换成内置 \(1/4\) 波介质片的圆波导和圆锥形喇叭，收端用矩形喇叭接收。
    \item 将矩形喇叭绕 \(z\) 轴旋转 \(180^\circ\)，测量旋转过程中电场强度的最小值 \(E_{\min}\) 和最大值 \(E_{\max}\)。
    \item 计算圆极化波的椭圆度 \(e = E_{\min}/E_{\max}\)，将数据填入表4-1。
\end{enumerate}

\subsection{圆极化波反射、折射特性的测试}
\begin{enumerate}
    \item 将接收喇叭换成内置 \(1/4\) 波介质片的圆喇叭。
    \item 设置入射角为 \(20^\circ\)，分别选择导体板及介质板作为反射板。
    \item 在接收端介质片法线 \(n'\) 与发射端介质片法线 \(n\) 处于 \(n' \perp n\) 和 \(n' \parallel n\) 两种相对位置时，测量直射波 \(E_i\)、反射波 \(E_r\) 和折射波 \(E_t\) 的幅度。
    \item 计算导体板的反射比 \(E_r/E_i\) 和介质板的能量关系 \((E_r^2 + E_t^2)/E_i^2\)，将数据填入表4-2。
\end{enumerate}

\section{实验数据}

\begin{table}[H]
    \centering
    \caption{表4-1 圆喇叭辐射圆极化波时椭圆度测试数据表}
    \renewcommand{\arraystretch}{1.3}
    \begin{tabular}{|c|c|c|c|c|}
        \hline
        圆极化波类型 &  接收 \(E_{\min}\) (V)/转角 & 接收 \(E_{\max}\) (V)/转角 & 椭圆度 \(e\) \\
        \hline
        右旋波 &  0.54/111\(^\circ\) & 1.63/23\(^\circ\) & 0.331\\
        \hline
    \end{tabular}
\end{table}

\begin{table}[H]
    \centering
    \caption{表4-2 圆极化波反射、折射特性测试数据表（入射角 \(20^\circ\)）}
    \renewcommand{\arraystretch}{1.3}
    \begin{tabular}{|c|c|c|c|c|}
        \hline
        圆极化波辐射状态 & \multicolumn{2}{c|}{\(\alpha = +45^\circ\)} \\
        \hline
        接收位置 \(n'\) 与 \(n\) 相对位置 & \(n' \perp n\) & \(n' \parallel n\) \\
        \hline
        接收直射波 \(E_i\) (V) & 0.46 & 1.94\\
        \hline
        接收导体板反射波 \(E_r\) (V) & 2.03 & 0.90 \\
        \hline
        接收介质板反射波 \(E_r\) (V) & 1.13 & 0.70 \\
        \hline
        接收介质板折射波 \(E_t\) (V) & 0.30 & 1.02\\
        \hline
        导体板的功率关系 (\(E_r/E_i\)) & 4.41 & 0.464\\
        \hline
        介质板的功率关系 (\(E_r^2 + E_t^2)/E_i^2\) & 6.455 & 0.406\\
        \hline
    \end{tabular}
\end{table}

\section{实验结果整理及误差分析}

\subsection{椭圆度测试结果分析}
从表4-1可知，右旋圆极化波的椭圆度实测值为 \(e = 0.331\)，与理想圆极化波（\(e = 1\)）存在较大差距。椭圆度偏低说明两个正交分量的幅度不等或相位差偏离 \(90^\circ\)。主要原因包括：介质片转角未精确调至使两分量幅度相等的角度；介质片长度 \(L\) 与工作频率的匹配存在偏差；系统初始对准误差及环境反射干扰等。实验中可通过微调发射频率（调整FREQ键）进一步优化椭圆度。

\subsection{反射、折射特性分析}
表4-2给出了入射角 \(20^\circ\) 时圆极化波的反射与折射测试结果：

\begin{itemize}
    \item 对于导体板：当 \(n' \perp n\) 时，反射波幅度（2.03 V）远大于直射波幅度（0.46 V），反射比 \(E_r/E_i = 4.41\)；当 \(n' \parallel n\) 时，反射波幅度（0.90 V）小于直射波幅度（1.94 V），反射比 \(E_r/E_i = 0.464\)。该结果表明导体板对圆极化波的反射具有明显的极化选择性，且 \(n' \perp n\) 时更利于接收反射波。
    
    \item 对于介质板：当 \(n' \perp n\) 时，反射波（1.13 V）和折射波（0.30 V）的能量平方和与入射波能量平方之比 \((E_r^2 + E_t^2)/E_i^2 = 6.455\)；当 \(n' \parallel n\) 时，该比值为 0.406。理想情况下该比值应接近1，实测值偏离理论值的原因包括：介质板表面存在多次反射与折射的干涉效应；系统未完全校准导致测量误差；环境杂波干扰；以及介质板本身的损耗等。
\end{itemize}

\subsection{误差来源分析}
\begin{enumerate}
    \item \textbf{系统对准误差：} 收发喇叭轴线未完全重合，或介质片法线方向与极化方向存在偏差，导致测量结果偏离理论值。
    \item \textbf{介质片参数偏差：} 介质片的长度 \(L\) 与工作频率的匹配程度直接影响相位差，进而影响圆极化质量。
    \item \textbf{环境反射干扰：} 实验室内的金属物体、墙壁等对电磁波的反射会叠加到测量信号中，造成幅度和相位的测量误差。
    \item \textbf{读数与记录误差：} 示波器读数的人为判读误差，以及螺旋测微器或转台刻度读数误差。
\end{enumerate}

综上所述，本实验验证了圆极化波的形成、辐射、接收及反射折射的基本特性，实测数据与理论预期趋势一致，但受限于实验条件和测量精度，部分数值存在一定偏差。

\end{document}